\chapter{Li Bai: The Banished Immortal}

\section{Life and legend}

Li Bai \zht{李白} (701--762), also romanised Li Po, courtesy name Taibai \zht{太白}.

\begin{itemize}
  \item Probably born in Central Asia (near Suyab, in modern Kyrgyzstan) and raised in
        Sichuan; his family background is obscure, perhaps merchant.
  \item Never sat the examinations. As a young man he wandered, practised swordsmanship,
        studied Daoism, and sought patrons.
  \item In Chang'an the elder poet He Zhizhang \zht{賀知章} read his verse and called him a
        \textbf{``banished immortal''} (\zht{謫仙人}): a heavenly being exiled to earth for
        some misdeed. The label stuck.
  \item 742--744: served Emperor Xuanzong as a court poet in the Hanlin Academy; anecdotes
        tell of him writing verses while drunk. He was eased out of court.
  \item 744: met Du Fu in Luoyang; they travelled together briefly. Du Fu wrote many poems
        about Li Bai; Li Bai wrote only a few about Du Fu.
  \item During the An Lushan rebellion he joined the staff of Prince Yong, who was then
        declared a rebel by his brother the new emperor. Li Bai was imprisoned and
        sentenced to exile in Yelang (far southwest), but pardoned on the way (759).
  \item Died 762. Legend: drunk in a boat, he leaned out to embrace the moon's reflection and
        drowned.
\end{itemize}

\begin{important}[Li Bai versus Du Fu]
\begin{center}
\begin{tabular}{lll}
\toprule
 & Li Bai \zht{李白} & Du Fu \zht{杜甫} \\
\midrule
Epithet & Immortal of Poetry \zht{詩仙} & Sage of Poetry \zht{詩聖} \\
Thought & Daoist, free, transcendent & Confucian, engaged, moral \\
Preferred form & Old-style songs, \emph{yuefu}, quatrains & Regulated octave \\
Style & Spontaneous, hyperbolic, fluent & Dense, crafted, allusive \\
Persona & The drinker, the wanderer & The worried official, the family man \\
\bottomrule
\end{tabular}
\end{center}
The pair are the twin peaks of the High Tang and the subject of endless comparison essays
(``which is greater?''). A good answer resists the stereotypes: Li Bai can be deeply
melancholy, Du Fu can be playful.
\end{important}

\section{The moon, again: ``Quiet Night Thoughts''}

We translated \zht{靜夜思} in Chapter 1. In the context of Li Bai's work, note:
\begin{itemize}
  \item It is a \emph{yuefu}-style quatrain (title of a song type), not regulated verse:
        \zht{明月} repeats, and the tonal pattern is irregular.
  \item The moon is Li Bai's most personal image: a companion, a mirror of home, a
        heavenly being like himself. He was said to have been named after the ``Great
        White'' star (Venus) seen in his mother's dream.
  \item Homesickness (\zht{思鄉}) is the theme, but Li Bai never really went home; the
        poem became the anthem of every Chinese person away from home.
\end{itemize}

\section{``Drinking Alone under the Moon'' \zht{月下獨酌} (no.\ 1 of 4)}

\begin{center}
\begin{tabular}{p{0.3\textwidth}p{0.6\textwidth}}
\toprule
\zht{花間一壺酒，} & Among the flowers, a jug of wine; \\
\zht{獨酌無相親。} & I drink alone, no friend beside me. \\
\zht{舉杯邀明月，} & Raising my cup, I invite the bright moon, \\
\zht{對影成三人。} & and with my shadow we make three. \\
\zht{月既不解飲，} & But the moon does not know how to drink, \\
\zht{影徒隨我身。} & and the shadow only follows me about. \\
\zht{暫伴月將影，} & Still, for now, with moon and shadow for company, \\
\zht{行樂須及春。} & we must make merry while it is spring. \\
\zht{我歌月徘徊，} & I sing, and the moon wavers to and fro; \\
\zht{我舞影零亂。} & I dance, and my shadow sprawls and scatters. \\
\zht{醒時同交歡，} & While sober we share our pleasure; \\
\zht{醉後各分散。} & once drunk, we each go our own way. \\
\zht{永結無情遊，} & Let us bind ourselves for ever in a friendship free of feeling, \\
\zht{相期邈雲漢。} & and meet again far off in the River of Clouds. \\
\bottomrule
\end{tabular}
\end{center}

\subsection{Analysis}
\begin{itemize}
  \item \textbf{Comic invention.} Loneliness turned into a party of three. The joke is
        worked out with logical care (the moon can't drink; the shadow just copies me).
  \item \textbf{Parallel couplet} \zht{我歌月徘徊，我舞影零亂}: ``I sing / I dance'',
        ``moon / shadow'', ``wavers / scatters''. The moon ``wavers'' because the drunk
        poet's vision sways; or the moon is dancing too.
  \item \textbf{\zht{無情遊}} ``a friendship without feeling'': companions that have no
        human emotions, so they can never betray or leave. A Daoist ideal of freedom from
        attachment, and also a sad joke about having no human friends.
  \item \textbf{\zht{雲漢}} = the Milky Way: the poet will meet his friends in heaven,
        where the banished immortal belongs.
  \item \textbf{Form.} Old-style five-character verse (14 lines), rhyme changes mid-poem.
\end{itemize}

\section{``Bring in the Wine'' \zht{將進酒}}

A \emph{yuefu} title (``Will you drink?''), with irregular lines of 3, 5, 7 and 10
characters. The opening:

\begin{center}
\begin{tabular}{p{0.38\textwidth}p{0.52\textwidth}}
\toprule
\zht{君不見黃河之水天上來，} & Have you not seen the Yellow River's waters coming from the sky, \\
\zht{奔流到海不復回。} & rushing to the sea, never to return? \\
\zht{君不見高堂明鏡悲白髮，} & Have you not seen in the high hall's bright mirror the grief of white hair, \\
\zht{朝如青絲暮成雪。} & black silk at dawn, by evening turned to snow? \\
\zht{人生得意須盡歡，} & When life goes your way, take all the joy you can; \\
\zht{莫使金樽空對月。} & don't let the golden cup face the moon empty. \\
\zht{天生我材必有用，} & Heaven made my talents: they must be of use. \\
\zht{千金散盡還復來。} & Spend a thousand in gold, and it will all come back. \\
\bottomrule
\end{tabular}
\end{center}

\begin{important}[Themes of ``Bring in the Wine'']
\begin{itemize}
  \item \textbf{Carpe diem} (\zht{及時行樂}): time is a river that never flows back, a
        life passes in a single day (``black silk at dawn, snow by evening'' --- hyperbole).
  \item \textbf{Self-confidence and frustration}: ``Heaven made my talents, they must be
        used'': and yet he is not being used.
  \item \textbf{Defiance}: later the speaker tells his friends to pawn his ``five-flower
        horse'' and thousand-gold furs for wine, ``to dissolve with you the sorrow of ten
        thousand ages'' (\zht{與爾同銷萬古愁}).
  \item \textbf{Form}: the \zht{君不見} ``Have you not seen'' formula comes from
        \emph{yuefu}; the irregular lines and sudden shifts of scale (cosmic river $\to$
        a single hair) give the sense of drunken improvisation, which was in fact carefully
        crafted.
\end{itemize}
\end{important}

\section{Regulated quatrains}

Li Bai was a master of the seven-character quatrain.

\subsection{``Setting Out Early from Baidi'' \zht{早發白帝城}}

\begin{center}
\begin{tabular}{p{0.38\textwidth}p{0.52\textwidth}}
\toprule
\zht{朝辭白帝彩雲間，} & At dawn I left Baidi among the coloured clouds; \\
\zht{千里江陵一日還。} & a thousand \emph{li} to Jiangling, back in a single day. \\
\zht{兩岸猿聲啼不住，} & On both banks the gibbons' cries never ceased --- \\
\zht{輕舟已過萬重山。} & my light boat has already passed ten thousand hills. \\
\bottomrule
\end{tabular}
\end{center}

Traditionally dated to 759, when Li Bai, on his way to exile, received his pardon at
Baidi in the Yangzi gorges and turned back downstream. The tonal pattern (a
heptasyllabic quatrain whose first line rhymes):

\begin{center}
\begin{tikzpicture}[x=0.8cm,y=0.8cm, every node/.style={font=\small}]
\node[font=\scriptsize, gray] at (1,1) {1};
\node[font=\scriptsize, gray] at (2,1) {2};
\node[font=\scriptsize, gray] at (3,1) {3};
\node[font=\scriptsize, gray] at (4,1) {4};
\node[font=\scriptsize, gray] at (5,1) {5};
\node[font=\scriptsize, gray] at (6,1) {6};
\node[font=\scriptsize, gray] at (7,1) {7};
\filldraw[fill=white, draw=black, dashed] (0.55,-0.45) rectangle (1.45,0.45);
\node at (1,0) {\zht{朝}};
\filldraw[fill=white, draw=black] (1.55,-0.45) rectangle (2.45,0.45);
\node at (2,0) {\zht{辭}};
\filldraw[fill=gray!45, draw=black, dashed] (2.55,-0.45) rectangle (3.45,0.45);
\node at (3,0) {\zht{白}};
\filldraw[fill=gray!45, draw=black] (3.55,-0.45) rectangle (4.45,0.45);
\node at (4,0) {\zht{帝}};
\filldraw[fill=gray!45, draw=black] (4.55,-0.45) rectangle (5.45,0.45);
\node at (5,0) {\zht{彩}};
\filldraw[fill=white, draw=black] (5.55,-0.45) rectangle (6.45,0.45);
\node at (6,0) {\zht{雲}};
\filldraw[fill=white, draw=black] (6.55,-0.45) rectangle (7.45,0.45);
\node at (7,0) {\zht{間}};
\node[font=\scriptsize, blue!60!black, anchor=west] at (7.6,0) {rhyme};
\filldraw[fill=white, draw=red!80!black, line width=1.6pt] (0.55,-1.55) rectangle (1.45,-0.65);
\node at (1,-1.1) {\zht{千}};
\filldraw[fill=gray!45, draw=black] (1.55,-1.55) rectangle (2.45,-0.65);
\node at (2,-1.1) {\zht{里}};
\filldraw[fill=white, draw=black, dashed] (2.55,-1.55) rectangle (3.45,-0.65);
\node at (3,-1.1) {\zht{江}};
\filldraw[fill=white, draw=black] (3.55,-1.55) rectangle (4.45,-0.65);
\node at (4,-1.1) {\zht{陵}};
\filldraw[fill=gray!45, draw=black] (4.55,-1.55) rectangle (5.45,-0.65);
\node at (5,-1.1) {\zht{一}};
\filldraw[fill=gray!45, draw=black] (5.55,-1.55) rectangle (6.45,-0.65);
\node at (6,-1.1) {\zht{日}};
\filldraw[fill=white, draw=black] (6.55,-1.55) rectangle (7.45,-0.65);
\node at (7,-1.1) {\zht{還}};
\node[font=\scriptsize, blue!60!black, anchor=west] at (7.6,-1.1) {rhyme};
\filldraw[fill=gray!45, draw=black, dashed] (0.55,-2.65) rectangle (1.45,-1.75);
\node at (1,-2.2) {\zht{兩}};
\filldraw[fill=gray!45, draw=black] (1.55,-2.65) rectangle (2.45,-1.75);
\node at (2,-2.2) {\zht{岸}};
\filldraw[fill=white, draw=black, dashed] (2.55,-2.65) rectangle (3.45,-1.75);
\node at (3,-2.2) {\zht{猿}};
\filldraw[fill=white, draw=black] (3.55,-2.65) rectangle (4.45,-1.75);
\node at (4,-2.2) {\zht{聲}};
\filldraw[fill=white, draw=black] (4.55,-2.65) rectangle (5.45,-1.75);
\node at (5,-2.2) {\zht{啼}};
\filldraw[fill=gray!45, draw=black] (5.55,-2.65) rectangle (6.45,-1.75);
\node at (6,-2.2) {\zht{不}};
\filldraw[fill=gray!45, draw=black] (6.55,-2.65) rectangle (7.45,-1.75);
\node at (7,-2.2) {\zht{住}};
\filldraw[fill=white, draw=black, dashed] (0.55,-3.75) rectangle (1.45,-2.85);
\node at (1,-3.3) {\zht{輕}};
\filldraw[fill=white, draw=black] (1.55,-3.75) rectangle (2.45,-2.85);
\node at (2,-3.3) {\zht{舟}};
\filldraw[fill=gray!45, draw=black, dashed] (2.55,-3.75) rectangle (3.45,-2.85);
\node at (3,-3.3) {\zht{已}};
\filldraw[fill=gray!45, draw=black] (3.55,-3.75) rectangle (4.45,-2.85);
\node at (4,-3.3) {\zht{過}};
\filldraw[fill=gray!45, draw=black] (4.55,-3.75) rectangle (5.45,-2.85);
\node at (5,-3.3) {\zht{萬}};
\filldraw[fill=white, draw=black] (5.55,-3.75) rectangle (6.45,-2.85);
\node at (6,-3.3) {\zht{重}};
\filldraw[fill=white, draw=black] (6.55,-3.75) rectangle (7.45,-2.85);
\node at (7,-3.3) {\zht{山}};
\node[font=\scriptsize, blue!60!black, anchor=west] at (7.6,-3.3) {rhyme};
\filldraw[fill=white] (0.6,-4.55) rectangle (0.9,-4.25); \node[font=\scriptsize, anchor=west] at (0.95,-4.4) {level \zht{平}};
\filldraw[fill=gray!45] (4.1,-4.55) rectangle (4.4,-4.25); \node[font=\scriptsize, anchor=west] at (4.45,-4.4) {oblique \zht{仄}};
\end{tikzpicture}
\end{center}

\begin{itemize}
  \item The only departure from the ideal pattern is \zht{千} in position 1 of line 2 (a free
        position).
  \item Rhyme words \zht{間} \emph{jiān}, \zht{還} \emph{huán}, \zht{山} \emph{shān}:
        lines 1, 2 and 4.
  \item \textbf{Speed.} Numbers (thousand \emph{li}, one day, ten thousand hills) and the
        word \zht{已} ``already'' make the poem rush. The gibbons, whose cries
        conventionally make travellers weep, are here only a blur of sound behind the boat:
        the convention is overrun by joy.
\end{itemize}

\subsection{``At Yellow Crane Tower, Seeing Off Meng Haoran to Guangling''
\zht{黃鶴樓送孟浩然之廣陵}}

\begin{center}
\begin{tabular}{p{0.38\textwidth}p{0.52\textwidth}}
\toprule
\zht{故人西辭黃鶴樓，} & My old friend takes leave of Yellow Crane Tower in the west; \\
\zht{煙花三月下揚州。} & in the third month's mist and flowers he goes down to Yangzhou. \\
\zht{孤帆遠影碧空盡，} & The distant shape of a lone sail vanishes into the blue; \\
\zht{唯見長江天際流。} & all I see is the Long River flowing to the edge of the sky. \\
\bottomrule
\end{tabular}
\end{center}

\textbf{Analysis.} The feeling of parting is never named. The last two lines simply record
the speaker's gaze following the sail until it is gone, and continuing to stare at the
river: the length of the look \emph{is} the friendship. A classic case of ``feeling in
scene'' (\zht{情景交融}, the fusion of feeling and scene).

\section*{Summary}
\begin{itemize}
  \item Li Bai: ``banished immortal''; Daoist, wanderer, drinker; Hanlin court poet
        742--744; caught up in Prince Yong's rebellion; legend of drowning reaching for the
        moon.
  \item Li Bai \zht{詩仙} vs Du Fu \zht{詩聖}.
  \item ``Drinking Alone under the Moon'': moon, shadow and self make three; friendship free
        of feeling.
  \item ``Bring in the Wine'': \emph{yuefu}, carpe diem, irregular lines, hyperbole.
  \item Quatrains: speed and joy (``Setting Out Early from Baidi''), feeling in scene
        (``Yellow Crane Tower'').
\end{itemize}
